\documentclass[11pt]{article}
\usepackage[margin=1.1in]{geometry}
\usepackage{amsmath,amssymb,amsthm}
\usepackage{array}

\newtheorem{theorem}{Theorem}
\newtheorem{proposition}[theorem]{Proposition}
\newtheorem{lemma}[theorem]{Lemma}
\newtheorem{corollary}[theorem]{Corollary}
\newtheorem{conjecture}[theorem]{Conjecture}
\newtheorem{openlemma}[theorem]{Open Lemma}
\theoremstyle{definition}
\newtheorem{definition}[theorem]{Definition}
\theoremstyle{remark}
\newtheorem{remark}[theorem]{Remark}

\newcommand{\Z}{\mathbb{Z}}
\newcommand{\R}{\mathbb{R}}
\newcommand{\E}{\mathbb{E}}
\newcommand{\CV}{\mathrm{CV}}
\newcommand{\Var}{\mathrm{Var}}
\newcommand{\cov}{\mathrm{cov}}
\newcommand{\lam}{\lambda}
\newcommand{\gam}{\gamma}
\newcommand{\bet}{\beta}
\newcommand{\kap}{\kappa}

\title{Towards density $x^{1-\varepsilon}$ for the $3x+1$ problem:\\
a program with one open lemma}
\author{Martien de Jong\\[3pt]
{\normalsize with computational collaboration: Jengo (AI)}}
\date{July 2026}

\begin{document}
\maketitle

\begin{abstract}
We assemble, as a single conditional theorem, a program that would push the
Krasikov--Lagarias density exponent for the $3x+1$ problem to its ceiling.
\emph{If} one uniform-attenuation lemma holds --- the intra-triple variation
of the edge certificate is contracted by a factor $\kap \leq \kap_{\max} < 1$
at every $3$-adic depth and scale --- \emph{then} the tempering exponent
$\alpha_k \to 1$, hence $\gam_k \to 1$, hence
$\#\{ n \leq x : \text{orbit reaches } 1\} \geq x^{1-\varepsilon}$ for every
$\varepsilon > 0$. The chain has six links, of which five are proved and one is
the Open Lemma. Proved links: (1)~the tempering law makes each certificate an
exactly solvable roulette measure raised to a scalar power; (2)~a mass
conservation identity, $3W_0 + W_2 + W_8 = 3$ precisely because
$2^{\log_2 3} = 3$, forcing subcriticality $a + c \leq 1$; (3)~a balance
identity, up-flow $= \tfrac{3}{16} =$ down-flow, so the drift is a pure
second-order effect; (4)~a directional low-pass, $\min(x + c\mathbf 1) =
\min(x) + c$, making the branch channel a one-way fine-mode filter; (6)~an
edge-rate theorem, $d\gam/dq = 1/\ln(4/3) = 3.4761\ldots$, which transfers
$q \to 1$ to $\gam \to 1$ analytically. The Open Lemma is stated with maximal
precision as an LP-duality statement: at the K--L optimum, complementary
slackness pins the binding entries, and dual concentration implies a response
slope $\leq 1 - \lam_{\min}$ on intra-triple variation. We give an honest
account of what would falsify the program --- three ``creep'' scenarios in
which $\kap \to 1$ --- and why the fine-end saturation (Prop.~23) and
depth-flat attenuation (Prop.~35) data disfavor each. We close by contrasting
this density-exponent route with Tao's 2019 almost-all log-density theorem:
different quantities, disjoint methods. All statements are labeled proved,
verified, measured, or open.
\end{abstract}

\section{Introduction}
\label{sec:intro}

Let $T(n) = n/2$ for even $n$, $T(n) = (3n+1)/2$ for odd $n$, and let
$\pi_1(x) = \#\{ n \leq x : \text{the } T\text{-orbit of } n \text{ reaches } 1\}$.
The $3x+1$ conjecture asserts $\pi_1(x) = \lfloor x \rfloor$; Lagarias
\cite{Lag85} surveys the problem and Tao \cite{Tao19} proves the strongest
known metric result. The strongest known \emph{density} lower bounds come from
the difference-inequality method of Krasikov and Lagarias \cite{KL03}: a finite
feasibility certificate at $3$-adic depth $k$ yields
$\pi_1(x) \geq x^{\gam_k}$ for all large $x$. Krasikov and Lagarias obtained
$\gam = 0.84$; the ceiling of the method is $\gam = 1$, i.e.\ density
$x^{1-\varepsilon}$ for every $\varepsilon$.

Two companion papers develop the ingredients. In \cite{Temp26} we verify
certificates to depth $k = 20$ ($\gam = 0.9146$), identify the certificate
eigenvectors as \emph{tempered roulette measures}
($\mathrm{eigvec}_k = \mathrm{roulette}^{\alpha_k}$, $R^2 > 0.99$), prove the
edge-rate theorem $d\gam/dq = 1/\ln(4/3)$, and reduce the full-density
statement to a single drift-direction inequality. In \cite{Drift26} we isolate
the mechanism of that drift: mass conservation and a balance identity (both
forced by $2^{\log_2 3} = 3$), a one-line directional low-pass, and the
measured attenuation constant $\kap_\infty = 0.839 \pm 0.002$, mechanistically
identified as the rigidity of binding constraints in a linear program. The
present paper is the \emph{synthesis}: it states the full program as one
conditional theorem, marks exactly where each link is proved and where the
single remaining lemma is invoked, states that lemma with maximal precision,
and confronts it honestly with the ways it could fail.

The result is not a proof of the $3x+1$ conjecture, nor even of its density
form. It is a proof that the density form follows from one clean, local,
falsifiable analytic statement --- and a demonstration that every other link in
the chain, including the two that at first look like the hard analytic content
(the transfer from homogenization to density, and the subcriticality of the
heterogeneity cascade), is already unconditional.

\paragraph{Labels.} Every assertion carries one of four labels.
\emph{Proved}: complete mathematical proof (here or in a cited companion).
\emph{Verified}: finite exact or floating-point computation, reproducible from
stated parameters. \emph{Measured}: statistical fit or extrapolation from
computed data. \emph{Open}: neither proved nor finitely checkable.
Section~\ref{sec:status} collects the ledger.

\paragraph{Notation.}
$\bet := \log_2 3 = 1.58496\ldots$; $\alpha$ is reserved for the tempering
exponent. $\CV$ denotes coefficient of variation (standard deviation over
mean). For $j \geq 1$, $[3^j] := \{ m \bmod 3^j : m \equiv 2 \pmod 3\}$, so
$|[3^j]| = 3^{j-1}$. At the edge $(\lam, \gam) = (2, 1)$ we write
\begin{equation}
W_0 = \lam^{-2} = \tfrac14, \qquad
W_2 = \lam^{\bet-2} = \tfrac34, \qquad
W_8 = \lam^{\bet-1} = \tfrac32 .
\label{eq:weights}
\end{equation}

\section{The Krasikov--Lagarias edge system}
\label{sec:kl}

Fix $k \geq 2$ and $\lam \in (1, 2]$. The nontrivial-tree K--L system
$L_k(\lam)$ \cite{KL03} seeks $c : [3^k] \to \R_{>0}$ with, for each
$m \in [3^k]$,
\begin{align}
m \equiv 2 \ (\mathrm{mod}\ 9): &\quad
c(m) \leq \lam^{-2} c(4m) + \lam^{\bet-2}\, \bar c\!\left(\tfrac{4m-2}{3}\right),
\label{eq:L1}\\
m \equiv 5 \ (\mathrm{mod}\ 9): &\quad
c(m) \leq \lam^{-2} c(4m),
\label{eq:L2}\\
m \equiv 8 \ (\mathrm{mod}\ 9): &\quad
c(m) \leq \lam^{-2} c(4m) + \lam^{\bet-1}\, \bar c\!\left(\tfrac{2m-1}{3}\right),
\label{eq:L3}
\end{align}
where $4m$ is reduced mod $3^k$, branch targets mod $3^{k-1}$, and
\begin{equation}
\bar c(t) := \min\{ c(t),\ c(t + 3^{k-1}),\ c(t + 2 \cdot 3^{k-1})\}
\label{eq:min}
\end{equation}
is the minimum over the three refinements (the \emph{triple}) of $t \in
[3^{k-1}]$. Feasibility of $L_k(\lam)$ implies $\pi_1(x) \geq x^{\gam}$,
$\gam = \log_2 \lam$, for all large $x$ \cite{KL03}. We call $m \mapsto 4m$ the
\emph{transport} (\emph{up-}) channel and the two branch maps the \emph{down-}
channel; all three class maps are bijections of $[3^k]$ resp.\ onto
$[3^{k-1}]$. The right-hand sides define a monotone, positively homogeneous
(topical) operator $F_\lam$; feasibility is equivalent to its nonlinear Perron
value being $\geq 1$, and the critical $\lam_k$ (the \emph{edge}) is located by
bisection, the Perron eigenvector by power iteration.

The verified certificates \cite{Temp26} are
\[
(k, \gam_k, q_k) :\quad
(13, 0.8624, 0.9519),\ (15, 0.8805, 0.9594),\ (17, 0.8953, 0.9655),
\]
\[
(19, 0.9069, 0.9705),\ (20, 0.9146, 0.9723),
\]
where $q_k$ is the min-loss parameter of \eqref{eq:minloss} below. The exponent
$\gam_k$ climbs steadily but with shrinking increments; the program explains
why, and what would have to be shown to push it to $1$.

\section{The conditional theorem}
\label{sec:main}

We first name the missing statement, then state the theorem it closes. The
full precision (LP-duality form) is deferred to Section~\ref{sec:openlemma}.

\begin{openlemma}[Uniform attenuation; informal]\label{open:informal}
There is a constant $\kap_{\max} < 1$, independent of depth $k$ and scale $P$,
such that at the Perron fixed point of the critical system $L_k(\lam_k)$ the
minimum \eqref{eq:min} contracts intra-triple variation by at most
$\kap_{\max}$: writing $\kap(P; k)$ for the ratio of the standard deviation of
the triple-minimum's difference-increment to that of its members
(Definition~\ref{def:kappa}),
\[
\kap(P; k) \;\leq\; \kap_{\max} \;<\; 1
\qquad \text{for all } P \text{ and all } k .
\]
\end{openlemma}

\begin{theorem}[Conditional full-density bound]\label{thm:main}
Suppose Open Lemma~\ref{open:informal} holds. Then the tempering exponent
$\alpha_k \to 1$, equivalently the min-loss parameter $q_k \to 1$, equivalently
the density exponent $\gam_k \to 1$; consequently, for every $\varepsilon > 0$,
\[
\pi_1(x) \;=\; \#\{ n \leq x : \text{orbit reaches } 1 \}
\;\geq\; x^{1 - \varepsilon}
\qquad \text{for all sufficiently large } x .
\]
\end{theorem}

The rest of Section~\ref{sec:main} is the proof of Theorem~\ref{thm:main},
assembled from six links. For each link we mark \textbf{[proved]} or
\textbf{[Open Lemma]} in boldface at its head, so that the logical dependency
is legible at a glance. Links (1)--(4) and (6) are proved unconditionally; the
Open Lemma enters only at link (5), where it converts the proved structural
low-pass into a quantitative uniform contraction.

\subsection{Link (1): tempering law and roulette exact solvability
\textnormal{[proved solvability; measured law]}}
\label{sec:link1}

At the level of natural density the Syracuse map is exactly solvable. By the
Terras bijection \cite{Ter76}, prescribing the first $r$ cascade exponents
confines an odd $n$ to one residue class mod $2^{\,(\text{sum})+1}$, so the
$w_i$ are independent geometric$(1/2)$ in density. The stationary measure of the
induced residue dynamics mod $3^j$ --- the \emph{roulette measure} $\pi^{(j)}$
--- is computable in closed form by downward recursion with no eigenproblem,
because the base map $r \mapsto 3r + 1 \bmod 3^j$ depends only on
$r \bmod 3^{j-1}$ (Proposition~4 of \cite{Temp26}, proved). At depth two,
\begin{equation}
\pi(1,2,4,5,7,8) \;=\; \frac{(8,\ 16,\ 11,\ 4,\ 2,\ 22)}{63}
\label{eq:mod9}
\end{equation}
(Theorem~5 of \cite{Temp26}, proved). The certificate eigenvectors are pure
powers of this measure,
\begin{equation}
\mathrm{eigvec}_k = \mathrm{roulette}^{\,\alpha_k},
\qquad R^2 = 0.9927 \ldots 0.9977,
\label{eq:tempering}
\end{equation}
with $\alpha_k = 0.8024, 0.8291, 0.8509, 0.8682, 0.8785$ at
$k = 13, 15, 17, 19, 20$ (Conjecture~T of \cite{Temp26}, measured at five
depths; the $k = 20$ value was an out-of-sample prediction confirmed to four
decimals). The role of this link is structural: it certifies that the entire
heterogeneity of the certificate is captured by the scalar deficit
$1 - \alpha_k$, so that ``homogenization'' has a one-parameter meaning.
Crucially, the linearization of $L_k(2)$ (replacing the minimum by the triple
mean) is exactly the roulette transition operator, whose spectrum is
$\{1\} \cup \{|\mu| = 1/4\}$ (verified at $j = 4, 5, 6$; Theorem~21 of
\cite{NOTE}, Remark~6 of \cite{Temp26}): linear damping toward the roulette is
overwhelming, so surviving heterogeneity must be re-injected at each level by
the min-nonlinearity. This is the mechanistic frame for links (2)--(5).

\subsection{Link (2): mass conservation and subcriticality
\textnormal{[proved]}}
\label{sec:link2}

\begin{lemma}[Mass conservation at the edge; proved, Lemma~24 of \cite{NOTE}]
\label{lem:mass}
With the edge weights \eqref{eq:weights}, the row masses of
\eqref{eq:L1}--\eqref{eq:L3} are $W_0 + W_2 = 1$ (type $2$), $W_0 = \tfrac14$
(type $5$), $W_0 + W_8 = \tfrac74$ (type $8$), and their average is exactly $1$:
\begin{equation}
3 W_0 + W_2 + W_8 \;=\; \tfrac34 + 3 \cdot 2^{\bet-2}
\;=\; \tfrac34 + \tfrac94 \;=\; 3,
\qquad \text{precisely because } 2^{\bet} = 3 .
\label{eq:massconv}
\end{equation}
\end{lemma}

\begin{proof}
Direct evaluation; the final step is
$3 \cdot 2^{\bet-2} = 3 \cdot 2^{\bet}/4 = 9/4$ iff $2^{\bet} = 3$.
\end{proof}

The difference field of the eigenvector propagates between adjacent trit scales
via the exact \emph{offset algebra} \cite{Temp26,Drift26}: a pure trit offset
$\delta = d \cdot 3^P$ satisfies $4\delta = \delta + 3\delta$ (transport:
scales $P, P{+}1$) and $\tfrac{4\delta}{3} = \tfrac{\delta}{3} + \delta$
(branch: scales $P{-}1, P$). Averaging over classes yields the two-coefficient
closure (measured, fit error $\leq 1.2\%$ at $k = 13, 15, 17, 19$)
\begin{equation}
\CV_P \;=\; a\, \CV_{P-1} + c\, \CV_{P+1},
\label{eq:lattice}
\end{equation}
where $\CV_P(k)$ is the mean $\CV$ over sibling triples differing only in trit
$P$ ($P = 1$ fine end, $P = k-2$ top).

\begin{corollary}[Subcriticality; proved within \eqref{eq:lattice}, Cor.\ of
Lemma~24]\label{cor:subcrit}
Each unit of difference-field mass at scale $P$ is redistributed to
$\{P-1, P, P+1\}$ with total weight equal to the row mass of its class
(Lemma~\ref{lem:mass}); the minimum \eqref{eq:min} is $1$-Lipschitz and
contributes no expansion. Averaging and applying the triangle inequality bounds
the closure coefficients, after absorbing the self-term, by
$a + c \leq 1$.
\end{corollary}

Measured at $k = 20$: $a + c = 0.9955$. The point is conceptual: the identity
$2^{\bet} = 3$ --- the defining equation of the problem --- pins the average row
mass to exactly $1$, so the cascade sits \emph{on} the conservative line and can
never grow exponentially. The only remaining freedom is the \emph{direction} of
drift along the line.

\subsection{Link (3): the balance identity
\textnormal{[proved]}}
\label{sec:link3}

\begin{proposition}[Balance identity; proved, Prop.~32(i) of \cite{NOTE},
Prop.~4 of \cite{Drift26}]\label{prop:balance}
At the edge, the up-flow of offset magnitude (scale $P \to P+1$, transport) and
the down-flow (scale $P \to P-1$, branch) are exactly equal:
\[
\text{up-flow} = W_0 \cdot \tfrac34 = \tfrac{3}{16}
= \bar w \cdot \tfrac14 = \text{down-flow},
\qquad \bar w := \tfrac{W_2 + W_8}{3} = \tfrac34 .
\]
The zeroth-order scale drift of difference mass is exactly zero.
\end{proposition}

\begin{proof}
Transport $4\delta = \delta + 3\delta$ splits magnitude $4|\delta|$ into
$|\delta|$ (scale $P$) and $3|\delta|$ (scale $P+1$); the up-going fraction is
$3/4$ and the transport weight $W_0 = 1/4$, so up-flow $= \tfrac14 \cdot
\tfrac34 = \tfrac{3}{16}$. Branch $\tfrac{4\delta}{3} = \tfrac{\delta}{3} +
\delta$ splits magnitude $\tfrac43|\delta|$ into $\tfrac13|\delta|$ (scale
$P-1$) and $|\delta|$ (scale $P$); the down-going fraction is $1/4$. The mean
branch weight per class is $\bar w = (W_2 + W_8)/3 = \tfrac34$ by mass
conservation; so down-flow $= \tfrac34 \cdot \tfrac14 = \tfrac{3}{16}$.
\end{proof}

The consequence is sharp and is the reason link (5) is needed at all:
\emph{the drift $c > a$ cannot be a first-moment effect}. Whatever pushes
heterogeneity preferentially downward must act through the single structural
asymmetry that survives the balance --- and there is exactly one.

\subsection{Link (4): the directional low-pass
\textnormal{[proved core]}}
\label{sec:link4}

\begin{proposition}[Directional low-pass; proved, Prop.~34b of \cite{NOTE},
Prop.~7 of \cite{Drift26}]\label{prop:lowpass}
For $x \in \R^3$ and $c \in \R$, $\min(x + c\mathbf 1) = \min(x) + c$, and
$|\min(x) - \min(y)| \leq \max_i |x_i - y_i|$. Hence, decomposing any triple
perturbation into its triple-constant (coarse) and intra-triple-varying (fine)
components, the minimum transmits coarse modes with slope exactly $1$ and is
$1$-Lipschitz --- weakly contractive --- on fine modes. Since only the
down-channel passes through the minimum, the edge operator transmits all modes
freely upward, transmits coarse modes freely, and attenuates precisely the
intra-triple-varying modes downward.
\end{proposition}

\begin{proof}
Every component of $x$ shifts by $c$, so the least does; and $x_i \leq y_i +
\max_j|x_j - y_j|$ for all $i$, so taking minima and symmetrizing gives the
Lipschitz bound.
\end{proof}

This is the entire structural content of the drift: the down-flow of fine modes
is starved by a factor $\kap \leq 1$ (with $\kap < 1$ whenever the triples are
not comonotone in their increments), while the up-flow is free. In
\eqref{eq:lattice} this is exactly the statement $c > a$. Proposition
\ref{prop:lowpass} is what makes the sign of the drift structural; what it does
\emph{not} give by itself is a \emph{uniform} bound $\kap \leq \kap_{\max} < 1$
across all depths. Supplying that bound is the Open Lemma (link (5)). The
supporting data are strong. Define:

\begin{definition}[Attenuation constant]\label{def:kappa}
For sibling triples at scale $P$ of the depth-$k$ certificate, let $D_i$ be the
member difference-increments under one application of the edge-linearized
operator and $D_{\min}$ that of the triple minimum; set
$\kap(P; k) := \mathrm{std}(D_{\min}) / \mathrm{std}(D_{\mathrm{member}})$.
\end{definition}

Measured (Prop.~35 of \cite{NOTE}; \cite{Drift26}): $\kap < 0.95$ at every
scale and depth, and the top-aligned deep-scale value is \emph{flat} in $k$:
\begin{equation}
\kap_{\mathrm{deep}} = 0.841\ (k{=}13),\ 0.840\ (15),\ 0.837\ (17),\
0.838\ (19)
\;\Longrightarrow\;
\kap_\infty = 0.839 \pm 0.002 .
\label{eq:kappadeep}
\end{equation}
Uniformity in $P$ and $k$ is thus \emph{empirically pinned}; the Open Lemma is
the analytic ratification of \eqref{eq:kappadeep}.

\subsection{Link (5): uniform attenuation
\textnormal{[Open Lemma]}}
\label{sec:link5}

This is the sole conditional step. We invoke Open Lemma~\ref{open:informal}
(precise form: Open Lemma~\ref{open:precise}, Section~\ref{sec:openlemma}) to
pass from the proved one-sided low-pass of link (4) to the quantitative,
depth-uniform contraction $\kap \leq \kap_{\max} < 1$. Combined with links
(2)--(4), this fixes the sign \emph{and} a uniform lower bound on the size of
the drift:
\begin{equation}
c - a \;\geq\; \eta(\kap_{\max}) \;>\; 0
\qquad \text{uniformly in } k,
\label{eq:drift}
\end{equation}
where $\eta(\cdot) > 0$ on $[\,0, 1)$ is the bookkeeping factor converting
one-sided low-passing into a gap in the emergent recurrence (the bare-to-
effective map; see the caveat in Section~\ref{sec:openlemma}). On the
conservative line $a + c \leq 1$ (Corollary~\ref{cor:subcrit}) with
$c - a \geq \eta > 0$, the decaying root of \eqref{eq:lattice} obeys
\begin{equation}
\theta_k \;=\; \frac{a_k}{c_k} \;\leq\; \bar\theta \;<\; 1
\qquad \text{uniformly in } k .
\label{eq:theta}
\end{equation}
The measured $\theta$-series is $0.8248, 0.8360, 0.8438, 0.8480, 0.8488$ at
$k = 13, 15, 17, 19, 20$, converging to $\theta_\infty \approx 0.85$
(consistent with $\kap_\infty = 0.839$ to $1.2\%$; the residual is the
up-channel mixing correction). Fine-end saturation
(Proposition~\ref{prop:cv1}, invoked also in Section~\ref{sec:falsify}) supplies
the nonzero source that makes \eqref{eq:theta} nonvacuous.

\subsection{Link (6): the edge rate transfers $q \to 1$ to $\gam \to 1$
\textnormal{[proved]}}
\label{sec:link6}

\begin{proposition}[Min-loss identity; proved, Thm.~12 of \cite{NOTE},
Prop.~9 of \cite{Temp26}]\label{prop:minloss}
Let $\lam_k$ be critical and $c$ the Perron eigenvector, so
\eqref{eq:L1}--\eqref{eq:L3} hold with equality. With
$S = \sum_{[3^k]} c$, $S_{\min} = \sum_{[3^{k-1}]} \bar c$,
\begin{equation}
1 \;=\; \lam_k^{-2} + \frac{q_k}{3}\left(\lam_k^{\bet-2} + \lam_k^{\bet-1}\right),
\qquad q_k := \frac{3 S_{\min}}{S} \in (0, 1],
\label{eq:minloss}
\end{equation}
and $q_k = 1$ iff every refinement triple of $c$ is constant, in which case
$\lam = 2$.
\end{proposition}

The identity follows by summing the edge equalities using that all class maps
are bijections; $q_k \leq 1$ with equality iff all triples are flat
\cite{NOTE,Temp26}. Verified to eight decimals at $k = 5, \ldots, 11$. It says:
\emph{the entire gap between the K--L method and full density is the min-loss
$1 - q_k$}.

\begin{theorem}[Edge rate; proved, Thm.~19 of \cite{NOTE}, Thm.~11 of
\cite{Temp26}]\label{thm:edgerate}
Along \eqref{eq:minloss}, with $\gam = \log_2 \lam$,
\[
\left.\frac{d\gam}{dq}\right|_{(\lam, q) = (2, 1)}
\;=\; \frac{1}{\ln(4/3)} \;=\; 3.4761\ldots,
\qquad\text{so } 1 - \gam \sim 3.4761\,(1 - q) .
\]
\end{theorem}

\begin{proof}
With $F(\lam, q) = \lam^{-2} + \tfrac{q}{3}(\lam^{\bet-2} + \lam^{\bet-1}) - 1$,
at $(2, 1)$ one computes $F_q = \tfrac34$ and $F_\lam = \tfrac{3(\bet-2)}{8}$
(using $2^{\bet-2} = \tfrac34$, $2^{\bet-3} = \tfrac38$), whence
$d\lam/dq = -F_q/F_\lam = 2/(2 - \bet)$ and
$d\gam/dq = \tfrac{1}{\lam \ln 2} \cdot d\lam/dq = \tfrac{1}{(2-\bet)\ln 2}
= \tfrac{1}{\ln 4 - \ln 3} = \tfrac{1}{\ln(4/3)}$.
\end{proof}

The measured ratios
$(1 - \gam_k)/(3.4761\,(1 - q_k)) = 0.824, 0.847, 0.873, 0.908$ at
$k = 13, 15, 17, 19$ increase monotonically to $1$ (Corollary of
Theorem~\ref{thm:edgerate}, measured): the transfer from homogenization to
density exponent is \emph{analytic}. In particular the empirical constant
$0.698$ in $\gam_k = 1 - 0.698\,(1 - \alpha_k)$ is a finite-$k$ composite,
$0.698 = 3.4761 \times (1 - q_k)/\CV_{\mathrm{res}}(k)$, with numerator and
denominator each tending to their edge limits.

\subsection{Assembly}
\label{sec:assembly}

Chaining the six links:
\begin{equation}
\underset{\text{link (5), Open Lemma}}{\kap \leq \kap_{\max} < 1}
\;\overset{(2),(3),(4)}{\Longrightarrow}\;
\theta_k \leq \bar\theta < 1
\;\Longrightarrow\;
\CV_{\mathrm{top}}(k) \lesssim \CV_1 \cdot \theta_k^{\,k} \to 0
\;\overset{(1)}{\Longleftrightarrow}\;
q_k \to 1
\;\overset{(6)}{\Longrightarrow}\;
\gam_k \to 1 .
\label{eq:chain}
\end{equation}
The middle equivalence uses the min--mean linearization
$E[1 - \min/\mathrm{mean}] = c_1 \CV + c_2 \CV^2$ with $c_1 \in [1.19, 1.45]$
(Prop.~18 of \cite{NOTE}, measured), so $1 - q_k = \Theta(\CV_{\mathrm{top}}(k))$,
together with $1 - \alpha_k = \CV_{\mathrm{res}}(k)$ (measured). The conclusion
$\gam_k \to 1$ is exactly $\pi_1(x) \geq x^{1-\varepsilon}$ for every
$\varepsilon > 0$, completing the proof of Theorem~\ref{thm:main}. \qed

Every link except (5) is unconditional. The Open Lemma is the only place the
argument leaves proved ground, and it enters at a single, sharply localized
point: the quantitative uniformity of a contraction whose \emph{sign} is
already proved (link (4)) and whose \emph{value} is already measured and
depth-flat (\eqref{eq:kappadeep}).

\section{The Open Lemma, precisely}
\label{sec:openlemma}

We now state the Open Lemma in its maximal-precision form: as a statement about
the dual of the Krasikov--Lagarias linear program. The mechanism identification
of \cite{Drift26} (measured) is that the attenuation is not clipping-at-switches
but \emph{selection-weighted mean reversion}: the argmin member of each triple
systematically co-moves less than the triple average, whether or not the
selection changes hands ($P(\mathrm{switch}) \approx 0.66$, flat; the
antisymmetric correction $\pm 0.026$ is carried equally by switch and no-switch
events, Prop.~34 of \cite{NOTE}). The LP reading is exact: the min-selected
member is the \emph{binding} constraint of the row that reads it, and binding
constraints are pinned by complementary slackness while slack members float.

\begin{openlemma}[Uniform attenuation at the fixed point; LP-duality form]
\label{open:precise}
Let $c_k$ be the Perron eigenvector of the critical system $L_k(\lam_k)$,
regarded as the optimal solution of the K--L linear program at $\lam = \lam_k$,
and let $y_k \geq 0$ be an optimal dual solution. For $t \in [3^{k-1}]$ let
$b(t) \in \{1, 2, 3\}$ be the binding (argmin) index of the triple over $t$
in \eqref{eq:min}, and let $\lam_{\min}(t)$ be the dual pinning strength of the
binding entry, i.e.\ the sensitivity coefficient
\[
\lam_{\min}(t)
\;:=\; 1 \;-\; \frac{\partial\, D_{\min}(t)}{\partial\, \bar D(t)} ,
\]
the shortfall between unit response of a free member and the response of the
binding member to a coarse (triple-constant) perturbation $\bar D(t)$ of the
eigenvector. Then there exists $\lam_\star > 0$, independent of $t$, $P$, and
$k$, such that at the fixed point
\[
\lam_{\min}(t) \;\geq\; \lam_\star
\qquad \text{whenever the triple over } t \text{ is non-degenerate,}
\]
provided the intra-triple coefficient of variation is bounded below,
\begin{equation}
\inf_k \CV_1(k) \;>\; 0 .
\label{eq:cv1pos}
\end{equation}
Equivalently, the response slope of the down-channel on intra-triple variation
satisfies $\kap(P; k) \leq 1 - \lam_\star =: \kap_{\max} < 1$ uniformly.
\end{openlemma}

\paragraph{Why complementary slackness pins the binding entries.}
At the critical $\lam_k$ the certificate saturates \eqref{eq:L1}--\eqref{eq:L3};
the down-channel of a row of type $2$ or $8$ reads its target triple only
through the minimum \eqref{eq:min}, so the KKT stationarity of the program
attaches a strictly positive dual multiplier to the argmin entry of every such
triple and zero to the two slack entries (complementary slackness). Coarse
perturbations of the eigenvector --- the triple-constant modes of link (4) ---
move slack and binding entries alike with slope $1$ through the \emph{level}
term $\lam^{-2} c(4m)$ (transport), but the binding entry additionally feeds the
saturated down-constraint, whose multiplier resists its motion. The resulting
response deficit is exactly $\lam_{\min}(t)$; dual concentration --- the claim
that the multiplier mass does not thin to zero on any binding entry as
$k \to \infty$ --- is what makes $\lam_{\min}(t) \geq \lam_\star > 0$ uniform.

\paragraph{Why dual concentration implies slope $\leq 1 - \lam_{\min}$.}
Sensitivity analysis of the LP: perturb the right-hand data by a coarse field
$\bar D$. Slack members respond with slope $1$ (Proposition~\ref{prop:lowpass},
coarse channel). A binding member's response is reduced by its pinning strength,
so its slope is $1 - \lam_{\min}(t)$. Averaging the member responses gives the
triple-minimum response slope
$\partial D_{\min}/\partial \bar D = 1 - \lam_{\min}(t)$, hence
$\kap(P;k) \leq \max_t (1 - \lam_{\min}(t)) \leq 1 - \lam_\star$. Non-degenerate
triple gaps --- guaranteed by \eqref{eq:cv1pos} --- keep the binding set stable
under small perturbations, so the reduction is uniform rather than washing out
in the limit.

\paragraph{Where the specific arithmetic must enter.}
Steps ``slack $\Rightarrow$ slope $1$'' and ``binding $\Rightarrow$ slope
$1 - \lam_{\min}$'' are standard LP sensitivity facts once dual concentration is
available. The unproved content is the \emph{uniform lower bound}
$\lam_\star > 0$ on the dual pinning strength. This is where the identities of
links (2)--(3) must be used: mass conservation (Lemma~\ref{lem:mass}) normalizes
the total dual mass to a fixed budget, and the balance identity
(Proposition~\ref{prop:balance}) prevents that budget from draining to either
end of the scale hierarchy --- neither the fine end (where it would be killed by
the vanishing of $\CV_P$, but \eqref{eq:cv1pos} forbids the source vanishing)
nor the top (where it would leave interior triples slack). A proof of Open
Lemma~\ref{open:precise} is therefore a proof that the defining arithmetic
$2^{\bet} = 3$, which pins the mass and balance identities, also pins the dual
mass away from zero on binding entries.

\paragraph{The bookkeeping caveat.}
The recurrence \eqref{eq:lattice} is \emph{emergent}, not a per-application
transfer law: a single application of the edge operator is scale-scattering
(60--68\% of any mode's output energy lands on the finest scale, regardless of
input scale; measured, \cite{Drift26}), so \eqref{eq:lattice} describes the
Perron fixed point, not one step. Consequently the conversion
$\kap \leq \kap_{\max} \Rightarrow c - a \geq \eta > 0$ (the factor $\eta$ in
\eqref{eq:drift}) is itself a fixed-point statement, and the ``bare''
coefficients read off naively from the one-step weights \emph{invert} the drift
(measured negative result, \cite{Drift26}). This is precisely why the Open
Lemma is posed at the fixed point and its proof route is variational
(LP-duality) rather than iterative. We flag this honestly: Open
Lemma~\ref{open:precise} closes the chain modulo the bare-to-effective
bookkeeping, and a fully rigorous write-up must discharge both together.

\paragraph{Measured targets the proof must reproduce.}
$\kap_\infty = 0.839 \pm 0.002$; $\lam_{\mathrm{clip}} := -\cov(\bar D, R)/
\Var(\bar D)$ growing $0.085 \to 0.140$ with depth; the antisymmetric
correction $\pm 0.026$ carried equally by switch and no-switch events; and the
exact clipping decomposition $\kap^2 = \Var(\bar D)/V + 2\cov/V + \Var(R)/V$
with the entire attenuation in the covariance term ($2\cov/V = -0.170 \to
-0.257$) and $\Var(R)/V \leq 0.028$ (verified exact identity, \cite{Drift26}).

\section{What would falsify this}
\label{sec:falsify}

The program is falsifiable, and its weakest joint is precisely the Open Lemma:
everything reduces to $\kap \leq \kap_{\max} < 1$ uniformly. There is exactly
one way the conditional theorem can be vacuous: the attenuation \emph{creeps} to
$1$ as depth grows, $\kap(P; k) \to 1$, so that $\theta_k \to 1$ and the
$\gam_k$-increments, though positive, sum to a limit strictly below $1$. We
enumerate the concrete creep scenarios and confront each with data.

\paragraph{Creep A: triples comonotonize.}
If the three members of each refinement triple became asymptotically comonotone
in their increments, the minimum would track the mean and $\kap \to 1$. This is
the direct negation of link (4)'s nontriviality. It is disfavored by
\emph{fine-end saturation}:

\begin{proposition}[Fine-end saturation; measured, nine depths, Prop.~23 of
\cite{NOTE}]\label{prop:cv1}
The finest-level triple $\CV$ of the certificate satisfies
\[
\CV_1(k) \;=\; 0.5136 \;-\; 0.337\,(0.910)^k
\]
with residual $10^{-6}$ across $8 \leq k \leq 20$: geometric convergence to the
finite, strictly positive limit $\CV_1^\infty \approx 0.514$.
\end{proposition}

Comonotonization requires $\CV_1 \to 0$; the measured saturation at a nonzero
limit forbids it. The raw intra-triple variation on which the low-pass acts is
permanently supplied, and its limit ($0.514$) is not marginal but order unity.
For Creep~A to occur, the nine-depth fit --- residual $10^{-6}$, a clean
geometric law with a positive asymptote --- would have to be an artifact that
reverses beyond $k = 20$, turning the asymptote downward. No mechanism for such
a reversal is known, and the linear edge spectrum (damping $1/4$, injection at
each level; link (1)) predicts a floor, not a decay.

\paragraph{Creep B: attenuation drifts up with depth.}
Even with $\CV_1$ saturated, $\kap$ could in principle drift toward $1$ if the
selection-weighted mean reversion weakened at large $k$. This is disfavored by
the \emph{depth-flatness} of the deep-scale attenuation:

\begin{proposition}[Kappa uniformity; measured, four certificates, Prop.~35 of
\cite{NOTE}]\label{prop:uniform}
$\kap(P; k) < 0.95$ at every scale and depth; the top-aligned deep-scale value
is flat in $k$: $\kap_{\mathrm{deep}} = 0.841/0.840/0.837/0.838$ at
$k = 13/15/17/19$, giving $\kap_\infty = 0.839 \pm 0.002$.
\end{proposition}

This is the empirically decisive datum. Every \emph{other} measured quantity of
the program drifts with depth --- $\alpha_k$, $q_k$, $\theta_k$ all move
monotonically --- but the attenuation does not: it is a \emph{constant} of the
cascade, not a running coupling. If anything the deep-scale value drifts very
slightly \emph{downward} ($0.841 \to 0.838$), the wrong direction for Creep~B.
For Creep~B, $\kap_{\mathrm{deep}}$ would need to turn around and climb through
$0.95$ and onward to $1$; four certificates spanning
$3^{12}\!\to\!3^{18}$ classes show no such trend.

\paragraph{Creep C: the bare-to-effective map degenerates.}
The subtlest scenario: $\kap$ stays bounded below $1$, yet the emergent gap
$c - a$ (factor $\eta$ in \eqref{eq:drift}) shrinks to $0$ because the
fixed-point bookkeeping converting one-sided low-passing into a profile drift
degenerates. This is not excluded by $\kap$-data alone --- it is the honest soft
spot flagged in Section~\ref{sec:openlemma}. It is disfavored, but only
indirectly, by the \emph{measured $\theta$-series} itself: $\theta_k =
0.8248 \to 0.8488$ with increments shrinking by $\approx 0.6$ per step, a
textbook geometric approach to a limit $\approx 0.85$, comfortably below $1$
and consistent with $\kap_\infty = 0.839$ to $1.2\%$. A degenerating
bare-to-effective map would show $\theta_k \to 1$; the observed convergence is
to $0.85$. Still, we rate Creep~C the residual risk: the balance identity
(link (3)) guarantees the drift is second-order, and a second-order effect
could in principle vanish faster than measured without contradicting any proved
statement. Closing Creep~C is exactly the content of proving the bookkeeping
alongside the Open Lemma.

\paragraph{Summary of falsifiability.}
The framework makes a sharp $k = 21$ prediction \cite{Temp26,Drift26}:
$\gam \approx 0.918$, $\theta \approx 0.850$, $\kap_{\mathrm{deep}} = 0.839 \pm
0.002$, fine-end saturation rate $0.910$, and a tempering fit
$R^2 > 0.9977$. Any material deviation --- $\kap_{\mathrm{deep}}$ rising,
$\CV_1$ turning down, $\theta$ overshooting $0.86$ --- falsifies the program at
its stated form. The data at five depths disfavor Creeps~A and B decisively and
Creep~C indirectly; none is excluded with the force of a proof, which is exactly
what the Open Lemma would supply.

\section{Relation to Tao (2019)}
\label{sec:tao}

Tao \cite{Tao19} proves that almost all Collatz orbits attain almost bounded
values: for any function $f(x) \to \infty$, the set of $n$ whose orbit ever
drops below $f(n)$ has \emph{logarithmic} density $1$. This is the strongest
known metric result on the problem. It is worth being precise about how it
relates to the present program, because the two are easy to conflate and are in
fact orthogonal.

\paragraph{Different quantity.}
Tao controls a \emph{logarithmic-density} statement about \emph{almost all}
starting values --- an $\text{almost-everywhere}$ / measure result with an
exceptional set of log-density zero (of unspecified, possibly large, natural
density). The present program targets a \emph{natural-density exponent}: an
explicit lower bound $\pi_1(x) \geq x^{\gam}$ on the \emph{count} of convergent
integers, pushing $\gam$ from $0.84$ toward $1$. Neither statement implies the
other. Tao's ``almost all in log-density'' does not yield any power-law count
$x^{\gam}$ with $\gam$ near $1$ in natural density; conversely $\gam \to 1$
would give $\pi_1(x) \geq x^{1-\varepsilon}$ but says nothing about the
remaining $\leq x^{1-\varepsilon}$ integers, and in particular does not reach
the almost-all conclusion in Tao's sense. Both fall short of the conjecture
$\pi_1(x) = \lfloor x \rfloor$, from opposite directions.

\paragraph{Different method.}
Tao's proof is \emph{analytic and probabilistic}: he studies the Syracuse random
variable on $3$-adic-like cyclic groups, establishes quantitative equidistribution
(a Fourier / stability estimate for the Syracuse valuation distribution), and
runs a first-moment / weighted argument over a random starting value. The
present program is \emph{spectral and structural on a finite linear feasibility
system}: it verifies exact Krasikov--Lagarias certificates, identifies their
eigenvectors as tempered roulette measures, and reduces the exponent limit to an
LP-duality statement about a nonlinear Perron fixed point. The shared object ---
the geometric-$w$ / Syracuse valuation law and its residue dynamics mod $3^j$
(our roulette measure, link (1); Tao's Syracuse distribution) --- is the same
underlying arithmetic, but it is used at opposite ends: Tao averages it against
a random orbit to get a metric almost-all statement; we temper it to fit a
deterministic certificate and read off a density exponent.

\paragraph{Complementarity.}
The two are complementary rather than competitive. Tao's result says the
\emph{typical} orbit descends; ours (conditionally) would say a
\emph{near-full-density set} provably reaches $1$ with an explicit exponent. A
resolution of the full conjecture presumably needs the pointwise content that
neither reaches --- for us, that the countable integer thread cannot ride the
measure-zero survivor set (the ``ultra-compressible needle'' of \cite{NOTE},
Prop.~13); for the metric route, the upgrade from log-density to natural
density and from almost-all to all. The programs illuminate different faces of
the same wall.

\section{Honest status}
\label{sec:status}

\begin{center}
\begin{tabular}{p{0.22\textwidth} p{0.71\textwidth}}
\hline
\textbf{Proved} &
Link (1) roulette solvability + closed form mod $9$ \eqref{eq:mod9}
(\cite{Temp26});
link (2) mass conservation $3W_0 + W_2 + W_8 = 3$ and subcriticality
$a + c \leq 1$ (Lem.~\ref{lem:mass}, Cor.~\ref{cor:subcrit});
link (3) balance identity, up-flow $= \tfrac{3}{16} =$ down-flow
(Prop.~\ref{prop:balance});
link (4) directional low-pass $\min(x + c\mathbf 1) = \min(x) + c$,
$1$-Lipschitz on fine modes, down-channel only (Prop.~\ref{prop:lowpass});
link (6) min-loss identity (Prop.~\ref{prop:minloss}) and edge rate
$d\gam/dq = 1/\ln(4/3)$ (Thm.~\ref{thm:edgerate});
feasibility $\Rightarrow$ density bound (\cite{KL03}). \\
\hline
\textbf{Verified} &
Certificates $k = 13$--$20$, $\gam = 0.8624$--$0.9146$;
identity \eqref{eq:minloss} to $8$ decimals at $k = 5$--$11$;
linear edge spectrum $\{1\} \cup \{|\mu| = 1/4\}$ at $j = 4, 5, 6$;
pointwise transport identity (residual $2.8 \times 10^{-4}$);
clipping decomposition as an exact identity; Gaussian Monte Carlo mismatch;
Terras cylinder densities. \\
\hline
\textbf{Measured} &
Tempering law $\alpha_k \uparrow$, $R^2 > 0.99$ (link (1));
$1 - \alpha_k = \CV_{\mathrm{res}}$; constant $0.698$;
edge-rate ratios $0.824 \to 0.908$;
closure \eqref{eq:lattice}, $a + c = 0.9955$ at $k = 20$;
$\theta$-series $0.8248 \to 0.8488$;
$\kap < 0.95$ everywhere, $\kap_\infty = 0.839 \pm 0.002$, depth-flat
(Prop.~\ref{prop:uniform});
$\lam_{\mathrm{clip}} = 0.085$--$0.140$, antisymmetric $\pm 0.026$, not
switch-exclusive;
fine-end saturation $\CV_1 \to 0.514 > 0$ (Prop.~\ref{prop:cv1});
scale-scattering $60$--$68\%$ to digit $0$; one-step reconstruction inversion;
out-of-sample hits at $k = 19, 20$. \\
\hline
\textbf{Open} &
Open Lemma~\ref{open:precise}: uniform attenuation
$\kap \leq \kap_{\max} < 1$ at the fixed point (LP-duality: dual concentration
on binding entries $\Rightarrow$ response slope $\leq 1 - \lam_{\min}$), given
$\inf_k \CV_1 > 0$; plus the bare-to-effective bookkeeping converting one-sided
low-passing into $c - a > 0$. \\
\hline
\end{tabular}
\end{center}

\medskip

The shape of the situation is unusual and worth stating plainly. The best known
rigorous handle on the density of $3x+1$ convergence --- the Krasikov--Lagarias
certificate --- is, to $R^2 > 0.99$ at a billion classes, one exactly solvable
measure raised to one scalar power (link (1)). The distance of that power from
$1$ \emph{is} the distance of the density exponent from $1$, through an explicit
derivative $1/\ln(4/3)$ (link (6)). The defining constant of the problem,
$2^{\bet} = 3$, pins the certificate's heterogeneity cascade to the conservative
line (link (2)) and balances its zeroth-order scale flow exactly (link (3)); the
only asymmetry left is that the downward channel passes through a minimum, which
transmits coarse structure perfectly and attenuates fine structure (link (4));
and that attenuation is measured to be a depth-independent constant,
$\kap_\infty = 0.839 \pm 0.002$, mechanistically identified as the rigidity of
binding constraints in a linear program. Nothing in the chain is mysterious.
One variational inequality --- the Open Lemma, LP-duality form --- remains to be
written down with full rigor, and the negative results show it must be proved at
the fixed point, not by iteration.

\subsection*{Reproducibility}
All certificates, verification logs, and the scripts generating every number in
this paper are archived in the research dossier accompanying this work
(directories \texttt{certificates/}, \texttt{scripts/}, \texttt{findings/}), and
in the two companion papers \cite{Temp26,Drift26}.

\subsection*{Acknowledgment of method}
The computations, fits, and drafts were produced in an extended human--AI
collaboration; every proved statement above has a self-contained proof in the
text or in a cited companion, independent of that provenance.

\begin{thebibliography}{9}

\bibitem{KL03}
I.~Krasikov and J.~C.~Lagarias,
\emph{Bounds for the 3x+1 problem using difference inequalities},
Acta Arith. \textbf{109} (2003), no.~3, 237--258. arXiv:math/0205002.

\bibitem{Lag85}
J.~C.~Lagarias,
\emph{The 3x+1 problem and its generalizations},
Amer. Math. Monthly \textbf{92} (1985), 3--23.

\bibitem{Tao19}
T.~Tao,
\emph{Almost all orbits of the Collatz map attain almost bounded values},
Forum Math. Pi \textbf{10} (2022), e12. arXiv:1909.03562 (2019).

\bibitem{Ter76}
R.~Terras,
\emph{A stopping time problem on the positive integers},
Acta Arith. \textbf{30} (1976), 241--252.

\bibitem{Temp26}
M.~de Jong (with Jengo, AI),
\emph{Krasikov--Lagarias certificates for the $3x+1$ problem are tempered
roulette measures}, companion paper, July 2026.

\bibitem{Drift26}
M.~de Jong (with Jengo, AI),
\emph{The drift of the Collatz cascade: balance identities, the directional
low-pass, and the attenuation constant}, companion paper, July 2026.

\bibitem{NOTE}
M.~de Jong (with Jengo, AI),
\emph{Six results on the Collatz map in family/pair coordinates}, research note,
July 2026 (Theorems 12, 19, 21; Lemma 24; Propositions 18, 23, 31--35;
Conjecture T).

\end{thebibliography}

\end{document}
